\documentclass[11pt]{article}
\usepackage[margin=1in]{geometry}
\usepackage{amsmath,amssymb}
\usepackage{graphicx}
\usepackage{booktabs}
\usepackage{hyperref}
\usepackage{xcolor}
\usepackage{caption}
\usepackage{multicol}
\usepackage{enumitem}
\hypersetup{colorlinks=true, linkcolor=blue, urlcolor=blue, citecolor=blue}

\title{Evasion Without Answers: An Empirical Study of the \emph{Shadow} \\
Keyword's Win-Rate Premium in a Custom \emph{Magic: The Gathering} Cube}

\author{
  \textit{GIRY Aurélien} \\
  \small Independent cube design project
}
\date{August 2026}

\begin{document}
\maketitle

\begin{abstract}
We investigate whether the \emph{Shadow} keyword ability -- an evasion mechanic under
which affected creatures can be blocked only by other Shadow creatures -- produces a
disproportionate competitive advantage in a custom, 360-card ``old-border'' retro cube
for \emph{Magic: The Gathering}. Combining automated draft simulation
(CubeCobra) with automated head-to-head match simulation (the open-source \emph{Forge}
game engine), we generated 100 eight-player drafts and played out the resulting
2{,}800 round-robin matches. An initial diagnosis on a 10-draft pilot sample suggested
that Shadow's apparent strength might be an artifact of the game engine's own
under-blocking behavior (mean legal-blocking rate $\approx 17$--$21\%$) rather than a
property of the mechanic itself. Reconstructing true legal blocking opportunities from
raw match logs and correlating them against Shadow-deck win rate on the full
100-draft corpus overturns this hypothesis: the correlation collapses from
$r=+0.689$ ($n=10$) to $r \approx +0.07$--$0.10$ ($n=91$--$180$, all
$p>0.27$) and is never negative, which is the opposite of the signature a
blocking-artifact explanation would predict. Decks whose winning seat carries at
least one Shadow card win $57\%$ of the 100 drafts studied. Cross-tabulating win
rate by color count and Shadow presence isolates two additive, statistically
independent effects: mono-color consistency contributes roughly $+5$ to $+7$
percentage points, while Shadow contributes an additional $+21$ to $+22$ points
regardless of color count ($p<10^{-15}$ for all four pairwise comparisons). We
conclude that Shadow's advantage in this specific cube's card pool is intrinsic
to the mechanic rather than an artifact of simulated-opponent blocking behavior,
while explicitly noting the observational (non-causal) nature of the evidence and
the small amount of corroborating human playtest data currently available.
\end{abstract}

\section{Introduction}

Cube design -- the construction of a large, singleton card pool intended for repeated
limited-format drafting -- is typically iterated on through subjective playtesting and
community consensus rather than controlled measurement. Claims that a given card or
mechanic is ``too strong'' are common in cube design communities but are rarely
backed by match-level data, in part because generating enough games to say anything
statistically meaningful about a specific card is prohibitively time-consuming for
human playtesters.

This paper documents an attempt to close that gap for one specific, falsifiable
question that arose during the iterative construction of a 360-card retro-styled cube
(pre-2003 card frame, or officially reprinted in that frame by more recent products):
\emph{does the Shadow keyword confer a win-rate advantage large enough to warrant a
design change, and if so, is that advantage a property of the mechanic itself or an
artifact of the automated opponent used to generate match data?}

Shadow is a static evasion ability first introduced in \emph{Tempest} (1997): a
creature with Shadow ``can block or be blocked by only creatures with shadow.''
Unlike flying or other partial-evasion keywords, a deck with zero Shadow creatures of
its own has \emph{no legal way to block} a Shadow attacker under any circumstance,
independent of board state. This unconditional character makes Shadow a natural
candidate for producing an inflated win rate specifically against an AI opponent that
under-utilizes blocking in general, since a general blocking weakness would be
invisible against Shadow (there is nothing to fail to do) while being fully exposed
against ordinary evasion-free creatures. Distinguishing ``Shadow is intrinsically
strong'' from ``the simulator cannot punish Shadow's downside-free evasion'' is
therefore the central methodological problem this paper addresses.

\section{The Cube and Simulation Pipeline}

\subsection{Card pool}

The cube under study contains 360 non-land cards drafted from Magic sets released
through late 2003, supplemented by mechanically newer cards later reprinted in an
official ``retro'' or ``old-border'' card frame (e.g. \emph{Time Spiral Remastered},
\emph{Ravnica Remastered}, \emph{Modern Horizons 2/3}, \emph{Innistrad Remastered},
\emph{The Brothers' War: Retro Frame Artifacts}), so that every physical card in the
eventual paper cube shares a consistent pre-2003 visual frame regardless of the
card's original mechanical release date. The cube's intended paper-draft format
allocates a 20-card mana-fixing pool (ten historical pain lands and their
corresponding ten guildgates) to a separate secondary land-draft phase, distinct
from the 360-card main pool. \emph{The CubeCobra-simulated corpus analyzed in this
paper does not implement this separation}: CubeCobra's bot-drafting simulator has
no mechanism for a secondary land sub-draft, so all 380 cards -- fixing lands
included -- were drafted as ordinary pack contents in a single undifferentiated
pool. This is a genuine divergence between the simulated environment analyzed here
and the cube's intended real-world paper format, discussed further in \S4.2. The cube supports both traditional one-versus-one and free-for-all (FFA)
multiplayer play; it was not designed with either format as an explicit priority,
and its playtesting history to date happens to have started with FFA sessions for
circumstantial rather than design reasons. The 8-player, one-versus-one
round-robin format used to generate the simulation data analyzed here is
consequently a legitimate target format in its own right, not merely a
simulation-driven compromise against an FFA-first design -- though it is also, separately,
the more practical format to simulate at scale, since Forge does not currently support
multiplayer free-for-all simulation.

\subsection{Draft simulation}

Draft pools were generated using CubeCobra's bot-drafting simulator, which uses a
learned per-card ``Elo'' rating and pack-signal heuristics to approximate human
drafting behavior across an 8-seat pod. One hundred independent 8-player drafts were
generated, yielding 800 constructed decks.

\subsection{Match simulation}

Each of the 800 decks was exported to Forge's native \texttt{.dck} format. Forge (an
open-source, rules-complete \emph{Magic: The Gathering} engine with a scriptable AI
opponent) was used to play a full round-robin within each 8-player draft pod
(\(\binom{8}{2}=28\) single-game matches per draft, 2{,}800 matches total), invoked
via Forge's headless \texttt{sim} command-line mode. Every match was logged at full
verbosity (attack/block declarations, zone changes, stack resolutions, life totals).

\subsection{Win-detection correction}

An initial analysis pipeline determined match winners from three redundant log line
patterns (\texttt{Match Winner -}, \texttt{Game Outcome: \dots{} has won}, and
\texttt{\dots{} has won!}). Manual log inspection revealed that AI-vs-AI matches that
reach Forge's per-game time limit (120{,}000\,ms) can emit two contradictory
\texttt{Game Outcome: \dots{} has won} lines in immediate succession for a single
game. Because the original pipeline credited whichever line it encountered first, at
least one match result was misattributed to the losing seat rather than the winner.
This was corrected by switching to the single, unambiguous \texttt{Match Result:
\textless{}seat\textsubscript{1}\textgreater{}: \textless{}score\textsubscript{1}
\textgreater{} \textless{}seat\textsubscript{2}\textgreater{}: \textless{}score
\textsubscript{2}\textgreater{}} summary line, which appears exactly once per match.
Across the full 100-draft corpus, 6 anomalous matches across 5 drafts were identified
and correctly resolved this way (4 AI-timeout cases and one genuine simultaneous
double-loss requiring a Forge-adjudicated replay game).

\subsection{Legal blocking-rate reconstruction}

A na\"ive count of \texttt{didn't block} log lines over-states the AI's blocking
passivity, because it does not distinguish (a) turns on which the defending player
controlled no creature at all, and (b) attacks that were legally unblockable given the
defender's board (Shadow attackers facing a non-Shadow board; flying attackers facing
a board with no flying or reach creature) from (c) genuine refusals to block a
creature the defender could legally have blocked. We reconstructed per-player board
state directly from the log by tracking creature-resolution and graveyard/exile
zone-change events (card ownership is unambiguous in a singleton card pool) and
cross-referenced attacker keywords against Forge's own card database (rather than a
hand-curated list) to classify every non-block decision into one of these three
categories. The \emph{legal blocking rate} is then defined as
\[
  \text{block\_rate} = \frac{\text{real\_block}}{\text{real\_block} +
  \text{real\_nonblock}},
\]
i.e. blocks divided by blocks plus \emph{genuine} refusals, excluding both empty-board
turns and legally-unblockable attacks from the denominator.

\section{Results}

\subsection{Blocking rate}

Across the 100 drafts, the legal blocking rate averages $23.71\%$ (SD $=4.37$,
range $13.36$--$34.29\%$). This is markedly higher than a na\"ive count would suggest
(which, on early pilot data, understated the rate by roughly half), but confirms that
Forge's AI opponent under-utilizes blocking relative to a competent human player in
absolute terms.

\subsection{Shadow dominance frequency}

In 57 of the 100 drafts, the tournament-winning deck (highest match record in its
8-player pod) carries at least one Shadow creature. Figure~\ref{fig:rank} shows the
distribution of final standings across all 180 (draft, seat) pairs in which a deck
carries at least one Shadow card: the modal outcome for a Shadow-carrying deck is an
outright win, with a roughly monotonic decline in frequency at lower standings.

\begin{figure}[h]
  \centering
  \includegraphics[width=0.62\textwidth]{fig_rank_distribution.pdf}
  \caption{Final standing distribution for the 180 (draft, seat) pairs in which the
  seat's deck contains at least one Shadow card, across 100 eight-player drafts.}
  \label{fig:rank}
\end{figure}

\subsection{Blocking rate versus Shadow performance}

An initial 3-draft pilot, followed by a 10-draft extension, found a strongly positive
correlation ($r=+0.689$) between a draft's legal blocking rate and the win rate of its
most Shadow-concentrated deck -- the \emph{opposite} sign to what an
under-blocking-artifact explanation predicts (that hypothesis predicts $r<0$: less
blocking should mechanically favor an unblockable attacker). Recomputing the same
three cuts of the data on the full 100-draft corpus (Table~\ref{tab:corr}) collapses
the correlation toward zero under every definition of ``the Shadow deck.''

\begin{table}[h]
\centering
\caption{Pearson correlation between per-draft legal blocking rate and Shadow-deck
win rate, for three ways of defining the relevant deck(s), on the 10-draft pilot
versus the full 100-draft corpus.}
\label{tab:corr}
\small
\begin{tabular}{p{5.6cm}ccc}
\toprule
Definition of ``Shadow deck'' & $n$ (100-draft) & $r$ (10-draft pilot) & $r$, $p$
(100-draft) \\
\midrule
Single most Shadow-concentrated deck per draft & 99 & $+0.689$ & $+0.073$, $p=0.47$ \\
All seats carrying $\geq 1$ Shadow card (splash included) & 180 & $+0.152$ & $+0.081$, $p=0.28$ \\
Seats with a concentrated package ($\geq 2$ distinct Shadow cards) & 91 & $+0.384$ & $+0.099$, $p=0.35$ \\
\bottomrule
\end{tabular}
\end{table}

None of the three 100-draft correlations is statistically distinguishable from zero.
Two individual drafts illustrate the pattern directly: \emph{draft\_54} has the
highest legal blocking rate in the entire corpus ($34.29\%$) and is nonetheless won by
a Shadow-carrying deck; \emph{draft\_100} has the lowest blocking rate in the corpus
($13.36\%$) and is won by a deck with zero Shadow cards. If the AI's general
under-blocking were the primary driver of Shadow's success, both cases should run in
the opposite direction.

\subsection{Color count, Shadow, and win rate}

Table~\ref{tab:crosstab} pools match results across all 800 decks by color count
(mono- versus multi-color, decks whose archetype label did not resolve to an
identifiable color count are reported separately and excluded from the comparisons
below) crossed with the presence or absence of at least one Shadow card in the
deck's registered card list.

\begin{table}[h]
\centering
\caption{Pooled match win rate by color count and Shadow presence (2-proportion
$z$-test against the same color-count group without Shadow).}
\label{tab:crosstab}
\begin{tabular}{lrrrc}
\toprule
Group & Decks & W/M & Win rate & vs.\ no-Shadow, same color count \\
\midrule
Mono-color, without Shadow   & 250 & 842/1750  & $48.11\%$ & -- (baseline) \\
Mono-color, with Shadow      & 116 & 569/812   & $70.07\%$ & $z=10.40$, $p=2.6\times10^{-25}$ \\
Multi-color, without Shadow  & 345 & 1032/2415 & $42.73\%$ & -- (baseline) \\
Multi-color, with Shadow     & 64  & 284/448   & $63.39\%$ & $z=8.06$,  $p=7.7\times10^{-16}$ \\
\bottomrule
\end{tabular}
\end{table}

\begin{figure}[h]
  \centering
  \includegraphics[width=0.62\textwidth]{fig_mono_shadow_crosstab.pdf}
  \caption{Pooled win rate by color count and Shadow presence ($n=2{,}562$ mono-color
  and $n=2{,}863$ multi-color matches in total). Error bars are omitted for
  legibility; all four pairwise gaps in Table~\ref{tab:crosstab} and
  Table~\ref{tab:additive} are significant at $p<0.05$.}
  \label{fig:crosstab}
\end{figure}

Table~\ref{tab:additive} decomposes the same four cells into two candidate effects: a
mono-color consistency premium (comparing mono to multi-color \emph{within} the same
Shadow status) and a Shadow premium (comparing with- to without-Shadow \emph{within}
the same color-count class).

\begin{table}[h]
\centering
\caption{Decomposition into a mono-color effect and a Shadow effect.}
\label{tab:additive}
\begin{tabular}{lrc}
\toprule
Comparison & $\Delta$ (points) & significance \\
\midrule
Mono $-$ Multi, without Shadow (``mono premium'', baseline) & $+5.38$ & $p=5.7\times10^{-4}$ \\
Mono $-$ Multi, with Shadow (``mono premium'', with Shadow) & $+6.68$ & $p=1.5\times10^{-2}$ \\
With $-$ without Shadow, mono-color (``Shadow premium'', mono) & $+21.96$ & $p=2.6\times10^{-25}$ \\
With $-$ without Shadow, multi-color (``Shadow premium'', multi) & $+20.66$ & $p=7.7\times10^{-16}$ \\
\bottomrule
\end{tabular}
\end{table}

Both the mono-color premium and the Shadow premium are individually significant and
of comparable magnitude regardless of the level of the other factor, indicating the
two effects are close to independent. A purely additive model
($\text{mono}\_\text{noShadow} + \Delta_{\text{Shadow}} + \Delta_{\text{mono}}$)
over-predicts the mono-with-Shadow cell by roughly 4 percentage points
($74.15\%$ predicted vs.\ $70.07\%$ observed), consistent with mild ceiling
compression as combined win rate approaches its upper bound rather than a
qualitatively different interaction. In all four comparisons, the Shadow effect
($\approx+21$ points) is three to four times larger in magnitude than the color-count
effect ($\approx+5$--$7$ points).

\section{Discussion}

\subsection{Interpretation}

The central finding is a dissociation between two candidate explanations for
Shadow's apparent strength that were, prior to this analysis, difficult to
distinguish from qualitative play testing alone. An under-blocking-artifact
account predicts that Shadow's advantage should shrink as the simulated
opponent's legal blocking rate rises; the data show no such relationship at
$n=100$, and the $n=10$ correlation that initially suggested one is better
explained as small-sample noise given how completely it collapses on replication.
The color-count crosstab further shows that Shadow's win-rate premium is not a
proxy for the previously-hypothesized mana-consistency advantage of mono-color
decks: both effects are present, of different magnitude, and largely additive.

\subsection{Limitations}

The evidence remains observational. Shadow-carrying decks are not randomly
assigned; they are the product of a drafting process in which a human or bot
drafter chose to commit to the archetype, and such decks may differ from
non-Shadow decks along dimensions other than the presence of Shadow cards
themselves (e.g. overall card quality of the surrounding shell, curve, or
removal density) that were not controlled for here. The finding that Shadow's
premium is stable across the mono/multi-color split partially addresses this
concern by ruling out one specific confound, but does not rule out others.

Separately, Forge's AI opponent's general under-blocking behavior
($\bar{x}=23.71\%$ legal blocking rate) remains a genuine simulation-fidelity
concern independent of the Shadow question: a rate this low likely inflates the
absolute win rate of \emph{any} aggressive or evasive strategy relative to
human play, even where (as shown here) it does not differentially explain
Shadow's specific advantage over other mechanics. The single available human
playtest data point for this cube (a two-player match between a Boros aggro
deck and a Golgari midrange deck, won 2--1 by the aggro deck) is consistent
with, but far too small a sample to independently confirm, the direction of
this general concern.

Finally, all data reported here come from 1-versus-1 round-robin play. The cube
supports both 1v1 and free-for-all (FFA) multiplayer without a strict design
preference for either (\S2.1), so this is not a mismatch against some other
``intended'' format -- but several mechanics evaluated qualitatively during this
project's design process (symmetric life-total taxes, edict effects, board-wide
damage sources) are known to behave differently between 1v1 and FFA regardless,
and we make no claim that Shadow's measured premium transfers unchanged to FFA
play specifically.

A related and more consequential divergence concerns mana fixing, and it bears
specifically on the mono-color premium reported in \S3.4 rather than on the
Shadow finding. Because the simulated corpus drafted all 20 fixing lands as
ordinary pack contents (\S2.1) rather than through the cube's intended separate
land draft, a multi-color deck in this corpus had to spend some of its limited
picks on fixing lands that a mono-color deck could skip entirely in favor of an
additional spell. This creates a plausible drafting-mechanism confound distinct
from any genuine in-game mana-consistency advantage: part of the measured
$+5$--$7$ point mono-color premium (Table~\ref{tab:additive}) could reflect
multi-color decks being systematically diluted with lower-impact picks by the
drafting process itself, rather than (or in addition to) mono-color decks
benefiting from more reliable mana during play.

We tested this directly on the existing corpus rather than leaving it as an open
question. Restricting to the 403 strictly two-color decks (a subset of the
409-deck bi/tri-color group in Table~\ref{tab:crosstab}, which also includes 6
three-color decks), we counted how many of the cube's 20 dedicated fixing lands
each deck's registered card list actually contained and tested this count against
win rate (Table~\ref{tab:fixingtiers}).

\begin{table}[h]
\centering
\caption{Two-color decks only: pooled win rate by number of dedicated fixing
lands drafted (0 of 20, 1--2, or 3 or more). Pearson correlation across all 403
decks: $r=-0.003$, $p=0.96$.}
\label{tab:fixingtiers}
\begin{tabular}{lrrr}
\toprule
Fixing lands in deck & Decks & Matches (W/M) & Win rate \\
\midrule
0  & 30  & 90/210    & $42.86\%$ \\
1--2 & 357 & 1170/2499 & $46.82\%$ \\
3+ & 16  & 42/112    & $37.50\%$ \\
\bottomrule
\end{tabular}
\end{table}

We find no relationship between fixing-land count and win rate: the correlation
is indistinguishable from zero, and win rate does not decline monotonically with
more fixing lands drafted (the 1--2 tier outperforms the 0 tier, the opposite of
what the opportunity-cost hypothesis predicts). The 3+ tier's lower rate is based
on only 16 decks and should not be over-interpreted. This provides no support for
the pick-opportunity-cost explanation of the mono-color premium; a genuine
in-game mana-consistency advantage, or some other unmeasured factor, remains the
more likely explanation, though we did not test those alternatives directly.
Critically, this confound (supported or not) cannot explain the Shadow finding
itself: the Shadow premium in Table~\ref{tab:additive} is measured \emph{within}
each color-count stratum (mono-with- vs.\ mono-without-Shadow; multi-with- vs.\
multi-without-Shadow), so any general pick-opportunity-cost effect from mixed
land drafting would apply symmetrically to both terms of each comparison and
cancel out of the estimated Shadow effect regardless. Whether the intended
separate land draft remains worth its added setup complexity for a human table --
independent of this statistical question -- is a design decision this paper does
not attempt to settle.

\subsection{Design implications}

Taken together, the results support treating Shadow's competitive advantage in
this cube as a property of the mechanic and card pool rather than an artifact
of the specific automated opponent used to generate the data, and therefore as
a legitimate candidate for a deliberate design intervention (reducing Shadow
card density, adding dedicated answers, or removing the sub-theme) rather than
something to be dismissed as a simulation quirk.

\section{Conclusion}

Across 100 automated drafts and 2{,}800 simulated matches, a Shadow-carrying deck
wins its draft pod 57\% of the time, and carrying at least one Shadow card is
associated with a win-rate premium of roughly 21 points, independent of the
deck's color count and essentially uncorrelated with the simulated opponent's
legal blocking rate ($|r|<0.10$ under three different data cuts, none
significant). An initial small-sample analysis had suggested the opposite
conclusion was plausible; replicating the analysis at ten times the sample size
was sufficient to overturn it. We take this as a case study in why cube
balance claims -- like most claims resting on a handful of observed games --
benefit from exactly this kind of scale before being acted on.

\section*{Acknowledgments}
This analysis was conducted as part of an iterative, multi-assistant cube design
process; log-parsing and simulation-pipeline code was developed and progressively
corrected across sessions with several large language model assistants, including
review and cross-checking of intermediate results between them. All numerical
results in this paper were independently recomputed from the raw
\texttt{shadow\_seats\_detail.csv}, \texttt{shadow\_vs\_blocking.csv}, and
\texttt{mono\_vs\_shadow.csv} export files rather than taken from any single
assistant's prior summary.

\clearpage
\appendix
\section{Full card pool}
\label{app:cardpool}

The 360-card draftable pool at the time the 100-draft corpus analyzed in this paper was generated, listed by color/type category. One card added to the pool after data collection (\emph{Doubtless One}, White) is included below for completeness but does not appear in the per-card statistics reported in \S3--4, since it was never drafted into any of the 100 pools. The cube's intended paper format allocates a further 20-card pool (Table~\ref{tab:fixingpool}) to a separate secondary land draft; as noted in \S4.2, the simulated corpus analyzed in this paper does not implement that separation and drafted these 20 cards alongside the 360-card main pool instead.

\subsection*{White (53)}
\begin{multicols}{3}
\footnotesize
\begin{itemize}[leftmargin=1em,itemsep=0pt,parsep=0pt]
\item Ajani's Pridemate
\item Akroma's Vengeance
\item Armageddon
\item Arrest
\item Aura of Silence
\item Banishing Light
\item Beloved Chaplain
\item Benevolent Bodyguard
\item Blinding Angel
\item Condemn
\item Decree of Justice
\item Disenchant
\item Doubtless One
\item Enlightened Tutor
\item Eternal Dragon
\item Exile
\item Gather the Townsfolk
\item Glorious Anthem
\item Glory
\item Humble
\item Intrepid Hero
\item Knight of Dawn
\item Land Tax
\item Longbow Archer
\item Lunarch Veteran
\item Master Decoy
\item Mentor of the Meek
\item Mobilization
\item Mother of Runes
\item Order of Leitbur
\item Order of the White Shield
\item Orim's Prayer
\item Pacifism
\item Paladin en-Vec
\item Path to Exile
\item Savannah Lions
\item Seal of Cleansing
\item Serra Angel
\item Shelter
\item Soltari Champion
\item Soltari Monk
\item Soltari Trooper
\item Soul Warden
\item Stonehorn Dignitary
\item Story Circle
\item Swords to Plowshares
\item Thraben Inspector
\item True Believer
\item Unquestioned Authority
\item Weathered Wayfarer
\item Whipcorder
\item White Knight
\item Wrath of God
\end{itemize}
\end{multicols}

\subsection*{Blue (51)}
\begin{multicols}{3}
\footnotesize
\begin{itemize}[leftmargin=1em,itemsep=0pt,parsep=0pt]
\item Air Elemental
\item Aquamoeba
\item Boomerang
\item Braingeyser
\item Brainstorm
\item Capsize
\item Careful Study
\item Circular Logic
\item Clone
\item Cloud of Faeries
\item Cloudskate
\item Counterspell
\item Daze
\item Deep Analysis
\item Delver of Secrets
\item Equilibrium
\item Fact or Fiction
\item Fade Away
\item Forbid
\item Force Spike
\item Force of Will
\item Frantic Search
\item Great Whale
\item Gush
\item Impulse
\item Man-o'-War
\item Mana Breach
\item Mana Leak
\item Memory Lapse
\item Merfolk Looter
\item Mulldrifter
\item Ninja of the Deep Hours
\item Opposition
\item Opt
\item Peregrine Drake
\item Ponder
\item Propaganda
\item Psionic Blast
\item Repulse
\item Rhystic Study
\item Serendib Efreet
\item Snap
\item Snapcaster Mage
\item Stormscape Familiar
\item Stroke of Genius
\item Thieving Magpie
\item Thing in the Ice
\item Time Warp
\item Wall of Tears
\item Wash Out
\item Wonder
\end{itemize}
\end{multicols}

\subsection*{Black (57)}
\begin{multicols}{3}
\footnotesize
\begin{itemize}[leftmargin=1em,itemsep=0pt,parsep=0pt]
\item Accursed Marauder
\item Ashen Ghoul
\item Black Knight
\item Blood Artist
\item Bone Shredder
\item Cabal Archon
\item Carnophage
\item Carrion Feeder
\item Chainer's Edict
\item Coffin Purge
\item Crypt Rats
\item Dark Confidant
\item Dark Ritual
\item Darkness
\item Dauthi Horror
\item Dauthi Slayer
\item Diabolic Edict
\item Dismember
\item Duress
\item Exhume
\item Faceless Butcher
\item Fallen Angel
\item Feed the Swarm
\item Festering Goblin
\item Gisa's Bidding
\item Graveborn Muse
\item Gravedigger
\item Gray Merchant of Asphodel
\item Gurmag Angler
\item Hymn to Tourach
\item Hypnotic Specter
\item Innocent Blood
\item Knight of Stromgald
\item Lord of the Undead
\item Mesmeric Fiend
\item Mindslicer
\item Nantuko Shade
\item Necropotence
\item Order of the Ebon Hand
\item Phyrexian Arena
\item Phyrexian Plaguelord
\item Phyrexian Rager
\item Phyrexian Reclamation
\item Profane Tutor
\item Putrid Imp
\item Ravenous Rats
\item Rotlung Reanimator
\item Shriekmaw
\item Snuff Out
\item Terror
\item Tragic Slip
\item Twisted Abomination
\item Unearth
\item Village Rites
\item Withered Wretch
\item Zombie Infestation
\item Zulaport Cutthroat
\end{itemize}
\end{multicols}

\subsection*{Red (50)}
\begin{multicols}{3}
\footnotesize
\begin{itemize}[leftmargin=1em,itemsep=0pt,parsep=0pt]
\item Abrade
\item Anger
\item Anger of the Gods
\item Arc Lightning
\item Avalanche Riders
\item Ball Lightning
\item Blasphemous Act
\item Chaos Warp
\item Earthquake
\item Faithless Looting
\item Fervor
\item Fiery Temper
\item Fireblast
\item Firebolt
\item Flametongue Kavu
\item Fling
\item Goblin Bombardment
\item Goblin Grenade
\item Goblin King
\item Goblin Lackey
\item Goblin Matron
\item Goblin Patrol
\item Goblin Ringleader
\item Goblin War Drums
\item Goblin Warchief
\item Grim Lavamancer
\item Hammer of Bogardan
\item Incinerate
\item Kird Ape
\item Lightning Bolt
\item Mogg Conscripts
\item Mogg Fanatic
\item Monastery Swiftspear
\item Pillage
\item Pyroclasm
\item Raging Goblin
\item Rorix Bladewing
\item Shock
\item Siege-Gang Commander
\item Skizzik
\item Spark Spray
\item Spellshock
\item Squee, Goblin Nabob
\item Starstorm
\item Stone Rain
\item Sulfuric Vortex
\item Thermo-Alchemist
\item Two-Headed Dragon
\item Viashino Sandstalker
\item Young Pyromancer
\end{itemize}
\end{multicols}

\subsection*{Green (53)}
\begin{multicols}{3}
\footnotesize
\begin{itemize}[leftmargin=1em,itemsep=0pt,parsep=0pt]
\item Acidic Slime
\item Alpha Status
\item Ancestral Mask
\item Avatar of Might
\item Basking Rootwalla
\item Beast Within
\item Birds of Paradise
\item Blastoderm
\item Brawn
\item Call of the Herd
\item Caller of the Claw
\item Constant Mists
\item Elvish Champion
\item Elvish Spirit Guide
\item Fecundity
\item Fyndhorn Elves
\item Genesis
\item Giant Growth
\item Harrow
\item Krosan Tusker
\item Lhurgoyf
\item Llanowar Elves
\item Naturalize
\item Nimble Mongoose
\item Overgrowth
\item Overrun
\item Plow Under
\item Pouncing Jaguar
\item Priest of Titania
\item Quirion Ranger
\item Rampant Growth
\item Rancor
\item Ravenous Baloth
\item River Boa
\item Roar of the Wurm
\item Seedborn Muse
\item Spike Feeder
\item Sylvan Library
\item Symbiotic Beast
\item Thragtusk
\item Uktabi Orangutan
\item Verdant Force
\item Verduran Enchantress
\item Wall of Blossoms
\item Wellwisher
\item Werebear
\item Wild Growth
\item Wild Mongrel
\item Wirewood Herald
\item Wirewood Symbiote
\item Worldly Tutor
\item Yavimaya Elder
\item Young Wolf
\end{itemize}
\end{multicols}

\subsection*{Colorless (40)}
\begin{multicols}{3}
\footnotesize
\begin{itemize}[leftmargin=1em,itemsep=0pt,parsep=0pt]
\item Altar of Dementia
\item Ashnod's Altar
\item Basalt Monolith
\item Black Vise
\item Bottle Gnomes
\item Burnished Hart
\item Chief of the Foundry
\item Chimeric Idol
\item Contagion Clasp
\item Darksteel Juggernaut
\item Fellwar Stone
\item Foundry Inspector
\item Gilded Lotus
\item Grindstone
\item Hedron Archive
\item Helm of Awakening
\item Helm of Possession
\item Howling Mine
\item Juggernaut
\item Junk Diver
\item Karn, Silver Golem
\item Lodestone Golem
\item Lotus Bloom
\item Lotus Petal
\item Mazemind Tome
\item Mesmeric Orb
\item Mind Stone
\item Mirrorworks
\item Myr Battlesphere
\item Nevinyrral's Disk
\item Ornithopter
\item Solemn Simulacrum
\item Steel Overseer
\item The Rack
\item Thran Dynamo
\item Tormod's Crypt
\item Triskelion
\item Voltaic Key
\item Wurmcoil Engine
\item Zuran Orb
\end{itemize}
\end{multicols}

\subsection*{Multicolored (37)}
\begin{multicols}{3}
\footnotesize
\begin{itemize}[leftmargin=1em,itemsep=0pt,parsep=0pt]
\item Absorb
\item Armadillo Cloak
\item Baleful Strix
\item Bladestitched Skaab
\item Blazing Specter
\item Bloodtithe Harvester
\item Chrome Courier
\item Coiling Oracle
\item Consume Strength
\item Edgewalker
\item Fire // Ice
\item Fires of Yavimaya
\item Firestorm Hellkite
\item Fleshtaker
\item Gerrard's Verdict
\item Goblin Anarchomancer
\item Goblin Electromancer
\item Goblin Legionnaire
\item Goblin Trenches
\item Jungle Barrier
\item Lightning Helix
\item Llanowar Dead
\item Mirari's Wake
\item Mystic Snake
\item Orcish Lumberjack
\item Probe
\item Prophetic Bolt
\item Psychatog
\item Putrefy
\item Recoil
\item Sabertooth Nishoba
\item Spiritmonger
\item Sterling Grove
\item Tempest Drake
\item Temporal Spring
\item Terminate
\item Vindicate
\end{itemize}
\end{multicols}

\subsection*{Utility Lands (19)}
\begin{multicols}{3}
\footnotesize
\begin{itemize}[leftmargin=1em,itemsep=0pt,parsep=0pt]
\item Barren Moor
\item Cephalid Coliseum
\item Darksteel Citadel
\item Faerie Conclave
\item Forbidding Watchtower
\item Forgotten Cave
\item Gemstone Mine
\item Ghitu Encampment
\item Goblin Burrows
\item High Market
\item Lonely Sandbar
\item Mishra's Factory
\item Rogue's Passage
\item Secluded Steppe
\item Spawning Pool
\item Stalking Stones
\item Tranquil Thicket
\item Treetop Village
\item Wasteland
\end{itemize}
\end{multicols}

\subsection*{Separately-drafted fixing pool (20)}
\label{tab:fixingpool}
Ten historical ``pain lands'' (one per ally/enemy color pair) and their corresponding ten guildgates, distributed to drafters via a dedicated secondary land draft rather than mixed into the main packs.
\begin{multicols}{3}
\footnotesize
\begin{itemize}[leftmargin=1em,itemsep=0pt,parsep=0pt]
\item Adarkar Wastes
\item Azorius Guildgate
\item Battlefield Forge
\item Boros Guildgate
\item Brushland
\item Caves of Koilos
\item Dimir Guildgate
\item Golgari Guildgate
\item Gruul Guildgate
\item Izzet Guildgate
\item Karplusan Forest
\item Llanowar Wastes
\item Orzhov Guildgate
\item Rakdos Guildgate
\item Selesnya Guildgate
\item Shivan Reef
\item Simic Guildgate
\item Sulfurous Springs
\item Underground River
\item Yavimaya Coast
\end{itemize}
\end{multicols}

\end{document}
